\documentclass[10pt,a4paper]{amsart}
\usepackage[margin=19mm]{geometry}
\usepackage{amsmath,amssymb,mathtools}
\usepackage[scr=rsfs,cal=euler]{mathalfa}
\usepackage{xfrac}
\usepackage[unicode,hidelinks,bookmarks=false]{hyperref}
\newcommand{\ie}{i.e.\ }
\newcommand{\cf}{cf.\ }
\newcommand{\resp}{respectively\ }
\newcommand{\Subsection}{\subsection}
\providecommand{\nobreakdash}{\nobreak\hbox{-}}
\setlength{\emergencystretch}{2em}
\multlinegap0pt
\usepackage{tikz}
\usepackage{amssymb}
\usepackage{braket}
\usepackage{mathtools}
\usepackage{xcolor}
\newtheorem{lemma}{Lemma}
\newtheorem{proposition}{Proposition}
\renewcommand{\L}{\length}
\def\R{\mathbb{R}}
\def\S{\mathbb{S}}
\def\length{\mathcal L}
\newcommand{\norm}[1]{{\left \Vert {#1}\right \Vert}}
\title{Homogeneous Sobolev gradient flow of the length functional}
\author{Philip Schrader, Glen Wheeler, Valentina Wheeler}
\date{Source: arXiv:2603.18504v2 (CC BY 4.0)}
\begin{document}
\maketitle

\noindent\emph{Source and licence.} Homogeneous Sobolev gradient flow of the length functional. Version arXiv:2603.18504v2 (CC BY 4.0), distributed by arXiv under the Creative Commons Attribution 4.0 International licence, http://creativecommons.org/licenses/by/4.0/, as recorded in the arXiv licence metadata for this exact version and shown on the abstract page https://arxiv.org/abs/2603.18504v2. Full licence notice: \texttt{licenses/arxiv-2603.18504-CC-BY-4.0.txt}.\par\smallskip

\noindent\textit{Editorial scope.} Counted unit: the article's proposition globalexistence (source0.tex lines 608--611). The article's complete original proof of it is delivered with it (source0.tex lines 612--693, one \texttt{proof} environment of 82 lines). No mathematics is written, repaired or re-derived: the statement and the proof are the article's own lines, in the article's own notation, and no retained line is substituted or re-flowed.  Retained uncounted as the source-local context the proof needs: lines 164--170; lines 388--394; lines 417--423; lines 491--495; lines 510--516; lines 560--562; lines 573--575; lines 722--728.  Nothing was omitted from any retained span, so the package carries no omission record.  No retained line contains a citation, so the wrapper carries no bibliography and no bibliography entry.  No retained line contains a figure environment, an image-inclusion command or drawing code, so no image is delivered and the package declares no asset dependency.  Editorial formatting only, in addition: an AMS \texttt{amsart} preamble that restates the article's own declarations and macros used by the retained lines, namely the environments \texttt{lemma}, \texttt{proposition} on separate counters, together with the standard packages those lines need.  The retained environments print in this document as Lemma 1, Proposition 1 in order of appearance, whereas the article prints the same items with the numbering of its own full document.  The complete native source of the article as distributed in the arXiv e-print for this exact version is bundled unchanged as \texttt{sources/196699/source0.tex}.
\begin{equation}\label{GF}\tag{GFa}
\partial_t X(u,t)=F_{\lambda,a}(X):= 
\begin{cases}
- \frac{\L(X)^{a-2}}{\lambda^2}\Big(X(u,t)+(X\ast_X \mathcal G^\lambda) (u,t)\Big), & \length(X)\neq 0 \\
0 & \length(X)=0
\end{cases}
\end{equation}

\begin{equation}\label{GF2}\tag{GF2}
\partial_t X(u,t)=F_{\lambda,2}(X)=
\begin{cases}
-\frac{1}{ \lambda^2} \left (X+X\ast_X\mathcal{G}^\lambda \right ) & \length(X)\neq 0 \\
0 & \length(X)=0
\end{cases}
\end{equation}

\begin{lemma}\label{bounds}
For any $\gamma\in W^{1,1}(\S,\R^2)$:
\begin{align}\label{Fest1}
\norm{F_\lambda(\gamma)}_{L^\infty}& \leq \frac{1}{2\lambda^2 }\length(\gamma), \quad \text{and} \\ 
\norm{(F_\lambda(\gamma))'}_{L^1}&\leq \frac{2}{\lambda^2}\length(\gamma).  \label{Fdest}
\end{align}
\end{lemma}

\begin{align}\label{Fa-est1}
\norm{F_{\lambda,a}(\gamma)}_{L^\infty}&\leq \frac{1}{2\lambda^2}\length(\gamma)^{a-1},\\
\label{Fa-est2}
\norm{(F_{\lambda,a}(\gamma))'}_{L^1}&\leq \frac{2}{\lambda^2}\length(\gamma)^{a-1}.
\end{align}

\begin{proposition}\label{existence}
For each $X_0\in W^{1,1}(\S,\R^2)$ with $\length(X_0)>0$, each $a\in\R$ and each $\lambda>0$, there exists $T_0>0$ and a unique
\[
X\in C^1([-T_0,T_0],W^{1,1}(\S,\R^2))
\]
such that $X(0)=X_0$ and $X$ satisfies \eqref{GF}.
\end{proposition}

\begin{lemma}\label{wellposedconstants}
For each $X_0\in W^{1,1}(\S,\R^2)$ with $\length(X_0)=0$, each $a\in\R$ and each $\lambda>0$, the unique $X\in C^1([0,\infty),W^{1,1}(\S,\R^2))$ satisfying \eqref{GF2} and $X(\cdot,0)= X_0$ is the stationary solution $X(\cdot, t)\equiv X_0$.
\end{lemma}

\begin{lemma}\label{apriori}
If $X\in C^1([-T_0,T_0],W^{1,1}(\S,\R^2))$ is a solution to \eqref{GF2}, then $\norm{X(\cdot, t)}_{L^\infty}$ and $\length(X(\cdot,t))$ are non-increasing. 
\end{lemma}

\begin{lemma} \label{lengthdecay}
Let $X$ be a solution to \eqref{GF2} with $X_0$ immersed. Abbreviate $\mathcal L (t)=\mathcal L(X(\cdot,t))$ and $\mathcal L_0=\mathcal L(0)$, then
\begin{align}\label{eq:ldecay}
\length(t)\leq \length_0 e^{-4t/(1+8\lambda^2)}
\,.
\end{align}
\end{lemma}

\begin{proposition}\label{globalexistence}
For every $X_0\in W^{1,1}(\S,\R^2)$ there exists a unique global solution
$X\in C^1([0,\infty),W^{1,1}(\S,\R^2))$ to \eqref{GF2} with $X(\cdot,0)=X_0$.
\end{proposition}

\begin{proof}

By Lemma \ref{wellposedconstants} it suffices to consider $\length(X_0)>0$.
Let
$
X\in C^1([0,T_{\max}),W^{1,1}(\S,\R^2))
$
be the positive-time-maximal solution given by Proposition~\ref{existence}.
By Lemma~\ref{bounds} and \ref{apriori},
\[
\norm{\partial_t X(\cdot,t)}_{W^{1,1}}
=\norm{F_\lambda(X(\cdot,t))}_{W^{1,1}}
\le \frac{5}{2\lambda^2}\length(X(\cdot,0))
\]
and so for any $t_1\leq t_2\leq T_{\max}$ we have
\[
\norm{X(\cdot,t_2)-X(\cdot,t_1)}_{W^{1,1}}\leq \Big\Vert \int_{t_1}^{t_2}\partial_t X(\cdot,t)dt\Big\Vert_{W^{1,1}}\leq |t_2-t_1|\frac{5}{2\lambda^2}\length(X(\cdot,0)).
\]
Therefore if $T_{\max}<\infty$ there exists $X_*\in W^{1,1}(\S,\R^2)$ such that
$X(\cdot,t)\to X_*$
in $W^{1,1}(\S,\R^2)$
as $t\uparrow T_{\max}.$

To apply Proposition~\ref{existence} with $X_*$ as initial value we require $\length(X_*)>0$.
Set $\ell(t):=\length(X(\cdot,t))=\norm{X'(\cdot,t)}_{L^1}$. 
Since $t\mapsto X'(\cdot,t)$ is $C^1$ as an $L^1$-valued map, we have
\[
|\ell(t_2)-\ell(t_1)|\leq \norm{X_u(\cdot,t_2)-X_u(\cdot,t_1)}_{L^1}\leq \int_{t_1}^{t_2}\norm{X_{ut}(\cdot, \tau)}_{L^1}d\tau
\]
hence $\ell$ is absolutely continuous and for a.e.\ $t<T_{\max}$,
\[
\ell'(t)\ge -\norm{X_{ut}(\cdot,t)}_{L^1}
=-\norm{F_\lambda(X(\cdot,t))'}_{L^1}
\ge -\frac{2}{\lambda^2}\ell(t),
\]
where we have used \eqref{Fa-est2}. Gr\"onwall's inequality then yields
\begin{equation}\label{lengthlowerbound}
\ell(t)\ge \ell(0)e^{-2t/\lambda^2}>0
\qquad\text{for all }t<T_{\max}.
\end{equation}
and by continuity of $\length$,
\[
\length(X_*)=\lim_{t\uparrow T_{\max}}\ell(t)\ge \ell(0)e^{-2T_{\max}/\lambda^2}>0.
\]
Now Proposition~\ref{existence} with $X_*$ as initial value gives a unique continuation of $X$ beyond time $T_{\max}$, contradicting the assumption of maximality. Therefore $T_{\max}=\infty$.
% Set $\ell(t):=\length(X(\cdot,t))=\norm{X'(\cdot,t)}_{L^1}$. 
% Since $t\mapsto X'(\cdot,t)$ is $C^1$ as an $L^1$-valued map, $\ell$ is absolutely continuous and for a.e.\ $t<T_{\max}$,
% \[
% \ell'(t)\ge -\norm{\partial_t X'(\cdot,t)}_{L^1}
% =-\norm{F_\lambda(X(\cdot,t))'}_{L^1}
% \ge -\frac{2}{\lambda^2}\ell(t),
% \]
% so Gr\"onwall's inequality yields
% \[
% \ell(t)\ge \ell(0)e^{-2t/\lambda^2}>0
% \qquad\text{for all }t<T_{\max}.
% \]
% If $T_{\max}<\infty$, then Lemmas~\ref{bounds}, \ref{lengthdecay}, and \ref{apriori} give
% \[
% \norm{\partial_t X(\cdot,t)}_{W^{1,1}}
% =\norm{F_\lambda(X(\cdot,t))}_{W^{1,1}}
% \le \frac{5}{2\lambda^2}\ell(t)
% \le \frac{5}{2\lambda^2}\ell(0)
% \]
% and
% \[
% \norm{X(\cdot,t)}_{W^{1,1}}
% \le \norm{X(\cdot,t)}_{L^\infty}+\ell(t)
% \le \norm{X_0}_{L^\infty}+\ell(0)
% \]
% for all $t<T_{\max}$. Hence $X(\cdot,t)$ is Cauchy in $W^{1,1}$ as $t\uparrow T_{\max}$, so there exists $X_*\in W^{1,1}(\S,\R^2)$ such that
% \[
% X(\cdot,t)\to X_*
% \qquad\text{in }W^{1,1}(\S,\R^2)
% \quad\text{as }t\uparrow T_{\max}.
% \]
% By continuity of $\length$,
% \[
% \length(X_*)=\lim_{t\uparrow T_{\max}}\ell(t)\ge \ell(0)e^{-2T_{\max}/\lambda^2}>0.
% \]
% Applying Proposition~\ref{existence} with $a=2$ to the initial datum $X_*$ gives a local solution to \eqref{GF2} beyond time $T_{\max}$, contradicting maximality. Therefore $T_{\max}=\infty$.
\end{proof}

\end{document}
